\documentclass[11pt]{article}
\usepackage[utf8]{inputenc}
\usepackage[T1]{fontenc}
\usepackage[margin=2.5cm]{geometry}
\usepackage{amsmath,amssymb,amsthm,mathtools}
\usepackage{enumitem}
\usepackage{xcolor}
\usepackage{hyperref}
\hypersetup{colorlinks=true,linkcolor=blue!60!black,urlcolor=blue!60!black}

\theoremstyle{plain}
\newtheorem{theorem}{Theorem}[section]
\newtheorem{proposition}[theorem]{Proposition}
\newtheorem{lemma}[theorem]{Lemma}
\newtheorem{corollary}[theorem]{Corollary}
\theoremstyle{definition}
\newtheorem{assumption}{Assumption}
\theoremstyle{remark}
\newtheorem{remark}[theorem]{Remark}

\newcommand{\R}{\mathbb{R}}
\newcommand{\N}{\mathbb{N}}
\newcommand{\E}{\mathbb{E}}
\newcommand{\KL}{\mathrm{KL}}
\newcommand{\Ham}{\mathrm{Ham}}
\DeclareMathOperator*{\argmin}{arg\,min}
\DeclareMathOperator*{\sumop}{\textstyle\sum}
\DeclareMathOperator{\supp}{supp}
\DeclareMathOperator{\dimm}{dim}

\title{\textbf{Detailed proof of the minimax lower bound\\
on the orbit space, in the $\phi$-chart}\\[0.3em]
\large (validation note for the orbital--MDL plan; multiplicative regime)}
\author{J.~Nemb\'e \\ \small (working note)}
\date{\today}

\begin{document}
\maketitle

\begin{abstract}
We give a complete proof of the minimax lower bound
$\inf_{\widehat f}\sup_{f^\star\in\Sigma_{\mathrm{inv}}(s,L)}\E\,h^2(P^\star,\widehat f)
\gtrsim n^{-2s/(2s+d)}$, $d=\dimm(E/G)$, for estimation of a $G$-invariant
density, working in the $\phi$-chart (multiplicative/twin regime). The argument
has three honest pillars: (a) a \emph{reduction} showing the orbit projection is
sufficient, so the problem is exactly density estimation on the $d$-dimensional
quotient; (b) a Varshamov--Gilbert hypothesis cube of $G$-invariant bumps with
controlled squared-Hellinger separation and Kullback--Leibler radius; (c) Fano's
inequality and a bias/complexity balance. The role of the $\phi$-chart is shown to
be \emph{transparent}: since squared Hellinger is invariant under the change of
variables, the bound transfers verbatim, and the heavy tails of the multiplicative
regime are irrelevant to the lower bound because the perturbations are compactly
supported. This validates the lower-bound side of the targeted contribution.
\end{abstract}

\section{Setting and the reduction to the quotient}

We use the notation of the orbital plan. $E$ carries a compact group $G$ acting by
isometries (in the $\phi$-chart $\phi(E)$); $\pi:E\to E/G$ is the quotient map,
$\mu$ a $G$-invariant reference measure, $\mu_G=\pi_*\mu$. Write $M:=E/G$ with the
infimum orbital metric $d_G$ and $d:=\dimm M=\dimm E-\dimm(\text{orbit})$. A
$G$-invariant density $f^\star=\tilde f^\star\circ\pi$ corresponds to a density
$\tilde f^\star$ on $(M,\mu_G)$. The squared Hellinger distance is
$h^2(P,Q)=\tfrac12\int(\sqrt{p}-\sqrt{q})^2\,d\mu$, and $H:=\sqrt{h^2}$ is the
Hellinger metric.

\begin{lemma}[Sufficiency of the orbit projection]\label{lem:suff}
Let $\mathcal F_{\mathrm{inv}}$ be a family of $G$-invariant densities on $E$ and
$X_1,\dots,X_n\stackrel{\text{iid}}{\sim}P_{f}$, $f\in\mathcal F_{\mathrm{inv}}$.
Then $\big(\pi(X_1),\dots,\pi(X_n)\big)$ is sufficient for $f$, are i.i.d.\ from
$\tilde P_{\tilde f}=\pi_*P_f$ on $(M,\mu_G)$, and
$h^2_E(P_f,P_g)=h^2_M(\tilde P_{\tilde f},\tilde P_{\tilde g})$. Consequently any
estimator $\widehat f(X_1,\dots,X_n)$ can be replaced, without increasing the
risk, by one depending only on $(\pi(X_i))_i$, and
\[
\inf_{\widehat f}\ \sup_{f\in\mathcal F_{\mathrm{inv}}}\E\,h^2_E(P_f,\widehat f)
\;=\;
\inf_{\widehat{\tilde f}}\ \sup_{\tilde f}\E\,h^2_M(\tilde P_{\tilde f},\widehat{\tilde f}).
\]
\end{lemma}

\begin{proof}
Invariance $f=\tilde f\circ\pi$ gives $f(X_i)=\tilde f(\pi(X_i))$, so the
likelihood factors through $\pi(X_i)$: sufficiency, and the pushforward of an
i.i.d.\ sample is i.i.d. For the Hellinger identity, invariant densities make the
affinity integrate to the quotient, $\int_E\sqrt{f g}\,d\mu=\int_M\sqrt{\tilde
f\tilde g}\,d\mu_G$. For the equality of minimax risks (which is all the lower
bound needs), note that under any $P_f$ with $f$ invariant, the conditional law of
$X_i$ given $\pi(X_i)=o$ is the (normalised) $G$-invariant measure on the orbit
$\pi^{-1}(o)$, which does \emph{not} depend on $f$: the within-orbit position is
\emph{ancillary} with $f$-free law. Hence the experiment $(X_{1:n};\{P_f\})$ is Le
Cam-equivalent to $(\pi(X_{1:n});\{\tilde P_{\tilde f}\})$ augmented by independent
$f$-free randomisation, and the two minimax risks coincide; squared Hellinger,
being invariant, is the same loss on both sides.
\end{proof}

Lemma~\ref{lem:suff} reduces the lower bound to density estimation of
$\tilde f^\star$ on the $d$-dimensional space $(M,\mu_G)$ from $n$ i.i.d.\ samples.
The $\phi$-chart enters only through the definition of smoothness (below); the
estimation problem itself is the standard one on $M$.

\section{Standing regularity}

\begin{assumption}[Regular chart on the quotient]\label{ass:reg}
There is a coordinate ball $U\subseteq M$ on which $(M,d_G,\mu_G)$ is Ahlfors
$d$-regular: $c_M^{-1}r^d\le \mu_G(B(x,r))\le c_M r^d$ for $x\in U$,
$0<r\le r_0$. Equivalently $U$ is bi-Lipschitz to a Euclidean ball in $\R^d$ with
$\mu_G$ comparable to Lebesgue measure. (Holds when the $G$-action is free with
smooth quotient; at non-free strata, $d$ is the dimension of the top stratum and
$U$ is taken inside it.)
\end{assumption}

We identify $U$ with a Euclidean cube $(0,1)^d$ and $\mu_G$ with Lebesgue measure
up to the constant $c_M$; this only affects constants below.

\begin{assumption}[H\"older class in the $\phi$-chart]\label{ass:holder}
$\Sigma(s,L)=\{\tilde f:\ \|\tilde f\|_{\mathcal C^s(M)}\le L\}$, the H\"older ball
of order $s>0$, smoothness measured in the $\phi$-chart, and
$\Sigma_{\mathrm{inv}}(s,L)=\{\tilde f\circ\pi:\tilde f\in\Sigma(s,L)\}$. Fix a
reference density $f_0\in\Sigma(s,L/2)$ with $0<c_0\le f_0\le C_1<\infty$ on $U$.
\end{assumption}

\section{The hypothesis cube ($G$-invariant bumps)}

Fix once and for all a $\mathcal C^\infty$ function $\psi:\R^d\to\R$ with
$\supp\psi\subseteq(0,1)^d$, $\int\psi=0$, $\|\psi\|_\infty\le1$,
$\|\psi\|_{\mathcal C^s}\le1$, and $\|\psi\|_{L^2}>0$.

For a bandwidth $h=1/m$ ($m\in\N$), partition $U$ into $N\asymp c_M\,h^{-d}$
disjoint sub-cubes $\{U_j\}_{j=1}^N$ of side $h$ and centres $x_j$. Define the
rescaled bumps and the hypothesis family
\begin{equation}\label{eq:bumps}
\psi_j(x):=L\,h^{s}\,\psi\!\Big(\frac{x-x_j}{h}\Big),\qquad
\tilde f_\omega := f_0+\sum_{j=1}^N\omega_j\,\psi_j,\qquad \omega\in\{0,1\}^N .
\end{equation}

\begin{lemma}[Admissibility]\label{lem:admiss}
For $h\le h_0$ (so that $L h^s\le c_0/2$), each $\tilde f_\omega$ is a probability
density on $M$ and $\tilde f_\omega\in\Sigma(s,c_1L)$ for a constant $c_1$
depending only on $\psi$.
\end{lemma}

\begin{proof}
$\int\psi_j=Lh^s\,h^d\int\psi=0$, so $\int\tilde f_\omega=\int f_0=1$; and
$|\sum_j\omega_j\psi_j|\le Lh^s\|\psi\|_\infty\le c_0/2$, so
$\tilde f_\omega\ge c_0/2>0$. The bumps have disjoint supports, and
$\|\psi_j\|_{\mathcal C^s}=Lh^s\cdot h^{-s}\|\psi\|_{\mathcal C^s}\le L$ (the
$s$-th difference quotient of a width-$h$, height-$Lh^s$ bump is $O(L)$); for
disjoint supports the global $\mathcal C^s$ seminorm of the sum equals the maximum
over $j$, so $\|\tilde f_\omega\|_{\mathcal C^s}\le\|f_0\|_{\mathcal C^s}+L\le c_1L$.
\end{proof}

\begin{lemma}[Varshamov--Gilbert {\cite[Lem.~2.9]{tsybakov}}]\label{lem:vg}
For $N\ge8$ there is $\Omega\subseteq\{0,1\}^N$ with $|\Omega|\ge2^{N/8}$ and
$\Ham(\omega,\omega')\ge N/8$ for all distinct $\omega,\omega'\in\Omega$. Fix such
an $\Omega$ and a base point $\omega^{(0)}\in\Omega$.
\end{lemma}

\section{Separation and information radius}

\begin{lemma}[Squared-Hellinger separation]\label{lem:sep}
For distinct $\omega,\omega'\in\Omega$,
\[
h^2(\tilde f_\omega,\tilde f_{\omega'})\ \ge\ \frac{1}{8C_1}\,\Ham(\omega,\omega')\,L^2h^{2s+d}\|\psi\|_{L^2}^2
\ \ge\ c_2\,L^2h^{2s},
\qquad c_2:=\frac{c_M\|\psi\|_{L^2}^2}{64\,C_1}.
\]
\end{lemma}

\begin{proof}
$(\sqrt a-\sqrt b)^2=(a-b)^2/(\sqrt a+\sqrt b)^2\ge(a-b)^2/(4C_1)$ since
$\tilde f_\omega,\tilde f_{\omega'}\in[c_0/2,C_1+c_0/2]\subseteq[0,2C_1]$ for
$h\le h_0$. Disjoint supports give
$\int(\tilde f_\omega-\tilde f_{\omega'})^2
=\sum_{j:\omega_j\ne\omega'_j}\int\psi_j^2
=\Ham(\omega,\omega')\,L^2h^{2s}\!\int\!\psi\big(\tfrac{\cdot-x_j}{h}\big)^2
=\Ham(\omega,\omega')\,L^2h^{2s+d}\|\psi\|_{L^2}^2$.
Hence $h^2=\tfrac12\int(\sqrt{}-\sqrt{})^2\ge\tfrac1{8C_1}\int(\cdot)^2$. Finally
$\Ham\ge N/8\asymp c_M h^{-d}/8$, so $h^2\ge c_2L^2h^{2s}$.
\end{proof}

\begin{lemma}[Kullback--Leibler radius]\label{lem:kl}
For every $\omega\in\Omega$, with $P_\omega^{\otimes n}$ the law of the sample,
\[
\KL\big(P_{\omega}^{\otimes n}\,\|\,P_{\omega^{(0)}}^{\otimes n}\big)
= n\,\KL(\tilde f_{\omega}\,\|\,\tilde f_{\omega^{(0)}})
\ \le\ \frac{2}{c_0}\,n\,N\,L^2h^{2s+d}\|\psi\|_{L^2}^2
\ \le\ c_3\,n\,L^2h^{2s}.
\]
\end{lemma}

\begin{proof}
$\KL(f\|g)\le\chi^2(f\|g)=\int(f-g)^2/g\le\frac{2}{c_0}\int(f-g)^2$, and
$\int(\tilde f_\omega-\tilde f_{\omega^{(0)}})^2=\Ham(\omega,\omega^{(0)})L^2h^{2s+d}\|\psi\|_{L^2}^2
\le N L^2h^{2s+d}\|\psi\|_{L^2}^2$. With $Nh^d\asymp c_M$ this is
$\le c_3 nL^2h^{2s}$ after the i.i.d.\ tensorisation
$\KL(P^{\otimes n}\|Q^{\otimes n})=n\KL(P\|Q)$.
\end{proof}

\section{Fano and the balance}

\begin{lemma}[Fano method {\cite[Thm.~2.7]{tsybakov}}]\label{lem:fano}
Let $\{P_\theta\}$, $(\Theta,H)$ a semimetric. If $\theta_0,\dots,\theta_{\!M}$
($M\ge2$) satisfy $H(\theta_j,\theta_k)\ge2\delta$ for $j\ne k$ and
$\frac1{M}\sum_{j=1}^{M}\KL(P_{\theta_j}\|P_{\theta_0})\le\alpha\log M$ with
$0<\alpha<1/8$, then
\[
\inf_{\widehat\theta}\max_{0\le j\le M}P_{\theta_j}\!\big(H(\widehat\theta,\theta_j)\ge\delta\big)
\ \ge\ \frac{\sqrt M}{1+\sqrt M}\Big(1-2\alpha-\sqrt{\tfrac{2\alpha}{\log M}}\Big)\ \ge\ c_\star>0 .
\]
\end{lemma}

\begin{theorem}[Minimax lower bound on the quotient, $\phi$-chart]\label{thm:lower}
Under Assumptions~\ref{ass:reg}--\ref{ass:holder} there are constants
$c,n_0>0$ such that for all $n\ge n_0$,
\[
\inf_{\widehat f}\ \sup_{f^\star\in\Sigma_{\mathrm{inv}}(s,L)}
\E\,h^2\big(P^\star,\widehat f\big)\ \ge\ c\,n^{-2s/(2s+d)} .
\]
\end{theorem}

\begin{proof}
By Lemma~\ref{lem:suff} it suffices to prove the bound on $(M,\mu_G)$. Take the
family $\{\tilde f_\omega:\omega\in\Omega\}$ of Lemmas~\ref{lem:admiss}--\ref{lem:vg};
all lie in $\Sigma(s,c_1L)$, hence (after relabelling $L\leftarrow L/c_1$) in
$\Sigma(s,L)$.

\emph{Choice of $h$.} Put $h=h_n$ with
\begin{equation}\label{eq:balance}
n\,L^2\,h^{2s+d}=c_4,\qquad\text{i.e.}\qquad h_n=\Big(\tfrac{c_4}{nL^2}\Big)^{1/(2s+d)},
\end{equation}
where $c_4$ is chosen below; then $N\asymp c_Mh_n^{-d}\to\infty$, so $h_n\le h_0$
and $N\ge8$ for $n\ge n_0$.

\emph{Separation.} By Lemma~\ref{lem:sep}, with $H=\sqrt{h^2}$,
$H(\tilde f_\omega,\tilde f_{\omega'})\ge\sqrt{c_2}\,L\,h_n^{s}=:2\delta_n$, i.e.
$\delta_n^2=\tfrac{c_2}{4}L^2h_n^{2s}$.

\emph{Information.} By Lemma~\ref{lem:kl},
$\frac1{|\Omega|}\sum_{\omega}\KL(P_\omega^{\otimes n}\|P_{\omega^{(0)}}^{\otimes n})
\le c_3\,nL^2h_n^{2s}=c_3\,c_4\,h_n^{-d}$ using \eqref{eq:balance}. Since
$\log|\Omega|\ge\frac{N}{8}\log2\ge c_5\,h_n^{-d}$, the Fano ratio is
\[
\frac{\frac1{|\Omega|}\sum_\omega\KL}{\log|\Omega|}\ \le\ \frac{c_3c_4}{c_5}\ =:\ \alpha .
\]
Choose $c_4$ small enough that $\alpha<1/8$.

\emph{Fano.} Lemma~\ref{lem:fano} gives
$\inf_{\widehat f}\max_\omega P_\omega(H(\widehat f,\tilde f_\omega)\ge\delta_n)\ge c_\star$.
By Markov,
\[
\inf_{\widehat f}\sup_{\tilde f\in\Sigma(s,L)}\E\,h^2(\widehat f,\tilde f)
\ \ge\ \delta_n^2\, c_\star
\ =\ \tfrac{c_2 c_\star}{4}L^2 h_n^{2s} .
\]
Finally, by \eqref{eq:balance}, $h_n^{2s}=(c_4/(nL^2))^{2s/(2s+d)}$, so
$L^2h_n^{2s}=c_4^{2s/(2s+d)}L^{2d/(2s+d)}\,n^{-2s/(2s+d)}$, whence the claimed
$c\,n^{-2s/(2s+d)}$ (the $L$-dependence is absorbed into $c$). Pulling back to
$E$ via Lemma~\ref{lem:suff} completes the proof.
\end{proof}

\section{Why the $\phi$-chart and the heavy tails are harmless here}

\begin{remark}[Chart transparency]
The construction lives on the quotient $M$ in the $\phi$-chart, where smoothness is
defined (Assumption~\ref{ass:holder}); but the loss $h^2$ is a change-of-variables
invariant, so the separations of Lemma~\ref{lem:sep} and the bound of
Theorem~\ref{thm:lower} are identical whether read in the $\phi$-chart or pulled
back to $E$ by $\phi^{-1}$. No Jacobian or branch issue enters the lower bound:
$\phi$ is fixed and bijective on the relevant compact ball.
\end{remark}

\begin{remark}[Heavy tails irrelevant to the lower bound]
The perturbations $\psi_j$ are supported in the fixed compact ball $U$ and the
hypotheses share the bounded reference $f_0$ there; tail behaviour of the
multiplicative regime ($\phi=\log$ on $(0,\infty)$, possibly heavy-tailed in the
original scale) plays no role. Heavy tails matter only for the \emph{upper} bound,
through the affinity/concentration step; the lower bound is a purely local,
compactly supported construction. Hence the matching rate
$n^{-2s/(2s+d)}$ is attained on multiplicative data without any tail assumption on
the lower-bound side.
\end{remark}

\begin{remark}[What this validates and what remains]
Theorem~\ref{thm:lower} discharges the lower-bound side of the targeted
contribution: combined with the orbital-MDL upper bound it yields the exact rate
on the quotient, including the multiplicative regime. The two remaining technical
points are both on the \emph{upper} side: (i) the twin Jackson estimate on a
possibly stratified quotient (orbifold strata), and (ii) the affinity/Hellinger
oracle constant under the twin-transported reference measure for heavy-tailed
positive data.
\end{remark}

\begin{thebibliography}{9}
\bibitem{tsybakov} A.~B.~Tsybakov, \emph{Introduction to Nonparametric
Estimation}, Springer, 2009 (Lemma~2.9, Theorem~2.7).
\bibitem{barroncover} A.~Barron, T.~Cover, ``Minimum complexity density
estimation,'' \emph{IEEE Trans.\ Inform.\ Theory} 37 (1991) 1034--1054.
\end{thebibliography}

\end{document}
